% ============================================================
%  07-chapitre1.tex  —  Cadre général et étude de l'existant
% ============================================================
\chapter{Cadre Général et Étude de l'Existant}

\begin{chapterintro}
  Ce premier chapitre établit le cadre général du projet SolarEase. Il présente
  l'organisme d'accueil, analyse le contexte du marché photovoltaïque, formule
  la problématique métier, étudie les solutions existantes et justifie la
  création d'une nouvelle plateforme unifiée.
\end{chapterintro}

% ──────────────────────────────────────────────────────────
\section{Présentation de l'Organisme d'Accueil}

\subsection{Informations générales}

\textbf{\organisme} est une jeune entreprise tunisienne fondée en
\textbf{2024} par \textbf{Marwen Zitouni} et \textbf{Iheb Haouech},
co-fondateurs et co-CEO. Implantée en Tunisie, elle se positionne sur trois
domaines technologiques complémentaires~: le \textbf{développement
logiciel}, les \textbf{solutions photovoltaïques} et les
\textbf{technologies intelligentes} (domotique et objets connectés).

L'entreprise adopte une approche \textit{full-stack} de l'innovation~:
de la conception logicielle jusqu'à l'intégration énergétique et la
recherche \& développement matérielle, en passant par le marketing
et la communication digitale. Cette transversalité lui permet de proposer
à ses clients des offres intégrées combinant logiciel, énergie et
domotique au sein d'un même projet.

\begin{table}[H]
  \caption{Fiche d'identité de \organisme}
  \label{tab:fiche-entreprise}
  \centering\small
  \begin{tabularx}{\textwidth}{|L{4.2cm}|X|}
    \hline
    \rowcolor{figblue}
    \textcolor{white}{\textbf{Élément}} &
    \textcolor{white}{\textbf{Description}} \\
    \hline
    Raison sociale & X\&V Electronics (XV Electronics) \\
    \hline
    Forme juridique & Société tunisienne \\
    \hline
    Année de création & 2024 \\
    \hline
    Co-fondateurs (CEO) & Marwen Zitouni \& Iheb Haouech \\
    \hline
    Secteurs d'activité &
    Développement logiciel \(\bullet\)~Solutions photovoltaïques
    \(\bullet\)~Domotique \& objets connectés \(\bullet\)~R\&D \\
    \hline
    Cible géographique &
    Tunisie, avec ambition régionale Afrique du Nord \\
    \hline
    Clientèle &
    Résidentielle, commerciale et industrielle (B2B et B2C) \\
    \hline
  \end{tabularx}
\end{table}

\subsection{Mission, vision et valeurs}

\noindent\textbf{Mission}~: créer des solutions technologiques innovantes
et adaptables pour les entreprises locales et internationales, en
combinant logiciel, énergie renouvelable et domotique au sein d'une
approche intégrée entre R\&D, production et marketing afin de générer
des résultats durables.

\noindent\textbf{Vision}~: devenir une référence en \textbf{Afrique du Nord}
dans l'intelligence digitale, l'énergie renouvelable et l'innovation
domotique en proposant des solutions sur mesure pour le marché tunisien
puis régional.

\noindent\textbf{Valeurs fondatrices}~:
\begin{itemize}
  \item \textbf{Innovation}~: prototypage de panneaux intelligents,
    intégration domotique et exploration de technologies émergentes
    (IA, IoT)~;
  \item \textbf{Qualité et fiabilité}~: choix des meilleures marques,
    tests rigoureux avant déploiement, suivi maintenance~;
  \item \textbf{Esprit d'équipe}~: collaboration étroite entre les pôles
    IT, énergie, R\&D et marketing~;
  \item \textbf{Satisfaction client}~: solutions adaptées aux besoins
    réels, accompagnement et support continu~;
  \item \textbf{Développement durable}~: utilisation responsable de
    l'énergie et promotion des sources renouvelables.
\end{itemize}

\subsection{Domaine d'activité}

\organisme{} structure son offre autour de quatre départements
opérationnels qui se complètent~:

\paragraph{a)~IT \& Développement logiciel.}
Conception et réalisation de plateformes web, d'applications mobiles,
de logiciels de gestion et de solutions SaaS. L'équipe développe
notamment des outils d'automatisation et de scalabilité tels que
\textit{LogiFlow} (suivi logistique), une plateforme tunisienne
inspirée d'Airbnb, des applications de livraison et des ERP sur mesure.
\textbf{SolarEase}, le projet objet de ce rapport, s'inscrit naturellement
dans cette activité IT en combinant développement logiciel et expertise
métier photovoltaïque.

\paragraph{b)~Solutions photovoltaïques.}
Études énergétiques, conception et installation de panneaux solaires
pour résidences, usines et entreprises. Le département assure également
la maintenance, l'optimisation des installations et l'intégration de
systèmes intelligents (supervision à distance, pilotage de la
consommation). L'engagement écologique se traduit par le choix des
meilleures marques et de solutions à fort rendement énergétique.

\paragraph{c)~R\&D et innovation.}
Prototypage de panneaux intelligents à \textbf{production nocturne},
innovations domotiques et tests de nouvelles technologies énergétiques.
C'est précisément dans ce département qu'est née la fonctionnalité
\textbf{Night Panel} intégrée à SolarEase, qui simule la production
résiduelle d'un panneau associé à un stockage par batterie.

\paragraph{d)~Marketing \& branding.}
Gestion des réseaux sociaux, identité visuelle, publicité digitale et
campagnes orientées conversion. L'équipe travaille en collaboration
étroite avec l'IT et la R\&D pour lancer rapidement les produits
commercialisables et soutenir la croissance de l'entreprise.

\subsection{Organisation interne}

L'entreprise adopte une structure simple et horizontale, adaptée à
une jeune société technologique. Les deux co-fondateurs assurent
conjointement la direction générale et coordonnent quatre départements
opérationnels, comme l'illustre la figure~\ref{fig:organigramme}.

\begin{figure}[H]
  \centering
  \begin{tikzpicture}[
      node distance=1.4cm and 2cm,
      box/.style={draw=isetblue, fill=isetblue!10, rounded corners=4pt,
                  minimum width=3.8cm, minimum height=1.1cm,
                  font=\small\bfseries, align=center, text width=3.6cm},
      topbox/.style={box, fill=isetblue, text=white,
                     minimum width=8cm, text width=7.6cm},
      link/.style={-latex, thick, isetblue},
      dlink/.style={-latex, dashed, gray, thin}
    ]
    \node[topbox] (dg)
      {Direction Générale\\
       {\footnotesize\mdseries Marwen Zitouni \& Iheb Haouech (co-CEO)}};

    \node[box, below=2cm of dg, xshift=-6cm] (it)
      {Pôle IT\\
       {\footnotesize Plateformes web,\\applis mobiles, SaaS}};
    \node[box, below=2cm of dg, xshift=-2cm] (energie)
      {Pôle Énergie\\
       {\footnotesize Installations PV,\\maintenance, optim.}};
    \node[box, below=2cm of dg, xshift=2cm] (rd)
      {Pôle R\&D\\
       {\footnotesize Night Panel,\\domotique, prototypes}};
    \node[box, below=2cm of dg, xshift=6cm] (mkt)
      {Pôle Marketing\\
       {\footnotesize Branding, social,\\publicité digitale}};

    \draw[link] (dg) -- (it);
    \draw[link] (dg) -- (energie);
    \draw[link] (dg) -- (rd);
    \draw[link] (dg) -- (mkt);

    \draw[dlink] (it.east) -- (energie.west)
      node[midway, above, font=\tiny\itshape]{intégration};
    \draw[dlink] (energie.east) -- (rd.west)
      node[midway, above, font=\tiny\itshape]{R\&D};
    \draw[dlink] (rd.east) -- (mkt.west)
      node[midway, above, font=\tiny\itshape]{go-to-market};
  \end{tikzpicture}
  \caption{Organigramme de \organisme}
  \label{fig:organigramme}
\end{figure}

\begin{table}[H]
  \caption{Rôles et responsabilités par pôle}
  \label{tab:roles-organisation}
  \centering\small
  \begin{tabularx}{\textwidth}{|L{3.8cm}|X|}
    \hline
    \rowcolor{figblue}
    \textcolor{white}{\textbf{Pôle}} &
    \textcolor{white}{\textbf{Responsabilités principales}} \\
    \hline
    Direction Générale &
    Pilotage stratégique, vision long terme, relations institutionnelles,
    validation des offres majeures et suivi financier global. Assurée
    conjointement par les deux co-fondateurs \\
    \hline
    Pôle IT \&\@ Développement &
    Conception et développement de plateformes web, applications mobiles,
    logiciels de gestion, ERP et solutions SaaS. Accent sur
    l'automatisation et la scalabilité \\
    \hline
    Pôle Énergie (Photovoltaïque) &
    Études énergétiques, conception, installation et maintenance de
    systèmes photovoltaïques pour clients résidentiels, commerciaux
    et industriels \\
    \hline
    Pôle R\&D &
    Prototypage de panneaux intelligents (production nocturne),
    innovations domotiques, tests de nouvelles technologies
    énergétiques, alimentation des autres pôles en innovations
    commercialisables \\
    \hline
    Pôle Marketing \&\@ Branding &
    Gestion des réseaux sociaux, identité visuelle, publicité
    digitale, campagnes de conversion et accompagnement du lancement
    des produits \\
    \hline
  \end{tabularx}
\end{table}

\noindent\textbf{Processus décisionnel}~: les décisions stratégiques sont
prises conjointement par les deux co-CEO en concertation avec les
responsables de pôles. Les choix techniques sont validés par les pôles
IT, Énergie ou R\&D selon leur nature, après revue et tests.

\noindent\textbf{Positionnement du projet SolarEase au sein de
\organisme}~: SolarEase est un projet \textbf{transverse} qui mobilise
trois pôles de l'entreprise~: le pôle \emph{IT \& Développement} pour
la conception et la réalisation logicielle de la plateforme, le pôle
\emph{Énergie} pour l'expertise métier (dimensionnement, choix des
équipements) et le pôle \emph{R\&D} pour la fonctionnalité Night Panel.
Le pôle \emph{Marketing} est sollicité pour la communication autour du
produit et son éventuelle commercialisation.

% ──────────────────────────────────────────────────────────
\section{Contexte et Enjeux du Projet}

\subsection{Le marché photovoltaïque en Tunisie}

Le plan solaire tunisien vise à porter la part des énergies renouvelables à
\textbf{35\%} de la production électrique nationale à l'horizon~2030. Dans ce
contexte, le nombre de sociétés d'installation photovoltaïque a fortement
progressé ces dernières années. La concurrence s'intensifie, et les sociétés
doivent impérativement professionnaliser leurs processus commerciaux et
techniques pour rester compétitives.

\subsection{Les enjeux opérationnels}

Les principaux défis auxquels font face les installateurs sont~:

\begin{table}[H]
  \caption{Enjeux opérationnels des installateurs photovoltaïques}
  \label{tab:enjeux}
  \centering
  \begin{tabularx}{\textwidth}{L{3.5cm} X L{3cm}}
    \toprule
    \rowcolor{figblue}
    \textcolor{white}{\textbf{Enjeu}} &
    \textcolor{white}{\textbf{Description}} &
    \textcolor{white}{\textbf{Impact}} \\
    \midrule
    Rapidité de traitement & Délai entre la demande client et la remise d'une offre & Perte de clients \\
    \addlinespace
    Précision technique & Erreurs de dimensionnement engendrant des sur-coûts & Litiges, insatisfaction \\
    \addlinespace
    Professionnalisme & Qualité variable des devis et propositions & Image et crédibilité \\
    \addlinespace
    Traçabilité & Difficulté à retrouver l'historique d'un dossier & Risque opérationnel \\
    \addlinespace
    Rentabilité & Absence de projection financière claire & Mauvaise argumentation \\
    \bottomrule
  \end{tabularx}
\end{table}

% ──────────────────────────────────────────────────────────
\section{Problématique et Objectifs}

\subsection{Problématique}

Les installateurs photovoltaïques en Tunisie utilisent aujourd'hui une multitude
d'outils hétérogènes et non intégrés (tableurs, simulateurs isolés, CRM génériques,
documents bureautiques). Cette fragmentation engendre des erreurs de dimensionnement,
des pertes de temps considérables, une absence de traçabilité et une image peu
professionnelle vis-à-vis du client final.

La problématique peut être formulée ainsi~:

\begin{center}
  \begin{tikzpicture}
    \node[draw=isetblue, fill=isetblue!8, rounded corners=6pt,
          inner sep=10pt, text width=0.8\textwidth, align=center,
          font=\itshape] {
      Comment concevoir et réaliser une plateforme web unifiée permettant
      à un installateur photovoltaïque de gérer l'intégralité de son cycle
      métier --- du premier contact client jusqu'au rapport de dimensionnement
      --- en s'appuyant sur des données fiables (PVGIS) et des recommandations
      intelligentes (IA)~?
    };
  \end{tikzpicture}
\end{center}

\subsection{Objectifs du projet}

Pour répondre à cette problématique, le projet SolarEase fixe les objectifs suivants~:

\begin{enumerate}
  \item Offrir une authentification sécurisée (JWT + OTP) avec gestion des rôles~;
  \item Centraliser la gestion des clients et des projets solaires dans un seul outil~;
  \item Automatiser le calcul de dimensionnement photovoltaïque à partir des données
    d'irradiation PVGIS~;
  \item Proposer un mode \textit{Night Panel} simulant le stockage par batterie~;
  \item Fournir une analyse financière complète (ROI, payback, cash-flow sur 25 ans)~;
  \item Générer automatiquement un rapport PDF professionnel et exploitable~;
  \item Intégrer une recommandation IA contextuelle via un LLM local (Ollama)~;
  \item Offrir un espace client dédié pour la consultation des projets et documents.
\end{enumerate}

% ──────────────────────────────────────────────────────────
\section{Analyse Critique de l'Existant}

\subsection{Outils actuellement utilisés}

L'analyse de l'environnement de travail révèle que les sociétés s'appuient
sur des outils non intégrés~:

\textbf{a)~Tableurs (Excel/Sheets)~:} La plupart des installateurs utilisent
des tableurs maison pour estimer la surface de toiture, le nombre de panneaux
et le coût d'installation. Ces outils sont sujets aux erreurs humaines, difficiles
à maintenir et non standardisés.

\textbf{b)~Simulateurs isolés~:} Des logiciels comme PVSyst ou PVGIS permettent
des calculs précis mais sont complexes à prendre en main et ne sont pas reliés
au processus commercial.

\textbf{c)~CRM généralistes~:} Les CRM de type Salesforce ou HubSpot offrent
une bonne gestion de la relation client mais ne disposent d'aucune logique
de dimensionnement solaire intégrée.

\textbf{d)~Documents bureautiques~:} Les propositions commerciales sont rédigées
manuellement sous Word, sans standardisation ni calcul automatique.

\subsection{Tableau comparatif des solutions existantes}

\begin{table}[H]
  \caption{Comparaison des solutions existantes}
  \label{tab:comparatif}
  \centering
  \small
  \setlength{\tabcolsep}{8pt}
  \renewcommand{\arraystretch}{1.12}
  \begin{tabularx}{\textwidth}{L{3.6cm} c c c c}
    \toprule
    \rowcolor{figblue}
    \multicolumn{1}{>{\centering\arraybackslash}p{3.6cm}}{\textcolor{white}{\textbf{Critère}}} &
    \multicolumn{1}{>{\centering\arraybackslash}p{1.6cm}}{\textcolor{white}{\textbf{Excel}}} &
    \multicolumn{1}{>{\centering\arraybackslash}p{1.6cm}}{\textcolor{white}{\textbf{PVSyst}}} &
    \multicolumn{1}{>{\centering\arraybackslash}p{2.0cm}}{\textcolor{white}{\textbf{CRM général}}} &
    \multicolumn{1}{>{\centering\arraybackslash}p{2.0cm}}{\textcolor{white}{\textbf{SolarEase}}} \\
    \midrule
    Dimensionnement & \textcolor{orange}{Partiel} & \textcolor{green!60!black}{Oui} & \textcolor{red!70!black}{Non} & \textcolor{green!60!black}{Oui} \\
    \addlinespace
    Gestion client & \textcolor{orange}{Partiel} & \textcolor{red!70!black}{Non} & \textcolor{green!60!black}{Oui} & \textcolor{green!60!black}{Oui} \\
    \addlinespace
    Analyse financière & \textcolor{orange}{Manuel} & \textcolor{orange}{Partiel} & \textcolor{red!70!black}{Non} & \textcolor{green!60!black}{Oui} \\
    \addlinespace
    Génération PDF & \textcolor{orange}{Partiel} & \textcolor{orange}{Partiel} & \textcolor{red!70!black}{Non} & \textcolor{green!60!black}{Oui} \\
    \addlinespace
    Aide IA & \textcolor{red!70!black}{Non} & \textcolor{red!70!black}{Non} & \textcolor{red!70!black}{Non} & \textcolor{green!60!black}{Oui} \\
    \addlinespace
    Données PVGIS & \textcolor{red!70!black}{Non} & \textcolor{green!60!black}{Oui} & \textcolor{red!70!black}{Non} & \textcolor{green!60!black}{Oui} \\
    \addlinespace
    Accès web & \textcolor{orange}{Partiel} & \textcolor{red!70!black}{Non} & \textcolor{green!60!black}{Oui} & \textcolor{green!60!black}{Oui} \\
    \addlinespace
    Coût & Faible & \textcolor{red!70!black}{Élevé} & \textcolor{red!70!black}{Élevé} & \textcolor{green!60!black}{Adapté} \\
    \bottomrule
  \end{tabularx}
  \vspace{6pt}
  \raggedright\footnotesize{Légende~: \textcolor{green!60!black}{Oui} = fonctionnalité complète; \textcolor{orange}{Partiel} = couverture partielle; \textcolor{red!70!black}{Non} = absence.}
\end{table}

\subsection{Conclusion de l'analyse}

L'étude de l'existant révèle une lacune claire~: \textbf{aucune solution
ne couvre l'intégralité du cycle métier} d'un installateur photovoltaïque
en une plateforme unique, accessible et abordable. SolarEase ambitionne
de combler ce vide en proposant une solution unifiée, sécurisée et intelligente.

% ──────────────────────────────────────────────────────────
\section{Positionnement et Apport de SolarEase}

SolarEase se distingue des solutions existantes par~:
\begin{itemize}
  \item \textbf{L'intégration complète}~: du premier contact client jusqu'au
    rapport PDF final, en passant par le dimensionnement et l'analyse financière~;
  \item \textbf{Le calcul basé sur des données réelles}~: connexion à l'API
    PVGIS de la Commission Européenne pour des données d'irradiance précises~;
  \item \textbf{L'innovation Night Panel}~: simulation et comparaison des
    panneaux avec stockage par batterie intégré~;
  \item \textbf{L'intelligence artificielle}~: recommandations contextualisées
    générées par un LLM local (Ollama / TinyLLaMA)~;
  \item \textbf{L'architecture moderne}~: microservices Spring Boot, frontend
    React TypeScript, déploiement Docker Compose.
\end{itemize}

\begin{center}
  \begin{tikzpicture}
    \node[draw=isetblue, fill=isetblue!10, rounded corners=8pt,
          inner sep=12pt, text width=0.75\textwidth, align=center,
          font=\itshape] {
      SolarEase fusionne les volets métier et technique dans un flux
      unique et automatisé~: de l'inscription jusqu'à l'édition
      d'un rapport PDF exploitable par le client final.
    };
  \end{tikzpicture}
\end{center}

% ──────────────────────────────────────────────────────────
\section{Méthodologie de Travail}

\subsection{Les méthodes agiles}

La méthode \textbf{Agile} est une approche de gestion de projet qui privilégie
la collaboration, la flexibilité et la livraison continue de valeur. Elle permet
à l'équipe de s'adapter en continu aux retours et aux évolutions des besoins,
sans figer un plan rigide dès le départ.

\subsection{La méthodologie Kanban}

\textbf{Kanban} est une méthode Agile à flux tiré, originellement développée
par Toyota pour optimiser les chaînes de production. Adaptée au développement
logiciel, elle repose sur la visualisation du travail sur un tableau divisé en
colonnes représentant les états successifs de chaque tâche.

Les principes fondamentaux de Kanban sont~:
\begin{itemize}
  \item \textbf{Visualisation du flux}~: toutes les tâches sont visibles
    sur le tableau (colonne À FAIRE, EN COURS, TERMINÉ)~;
  \item \textbf{Limitation du WIP}~: le nombre de tâches simultanées en
    cours est limité pour éviter la surcharge~;
  \item \textbf{Flux continu}~: contrairement à Scrum, il n'y a pas de
    cycles fixes~; les tâches avancent au rythme de la capacité réelle~;
  \item \textbf{Amélioration continue}~: le tableau permet d'identifier
    les goulots d'étranglement et d'y remédier rapidement.
\end{itemize}

Pour le projet SolarEase, le tableau Kanban a été géré sur \textbf{Jira}
(plateforme Atlassian). Le détail du Product Backlog, de l'équipe et de la
planification des sprints est présenté au chapitre suivant (section~2.3).

\section*{Conclusion}

Ce chapitre a permis de poser le cadre général du projet en présentant
l'organisme d'accueil, le contexte du marché photovoltaïque tunisien et les
enjeux opérationnels des installateurs. L'analyse comparative des solutions
existantes confirme l'absence d'un outil intégré couvrant l'intégralité du
cycle métier. Face à ce constat, SolarEase se positionne comme une plateforme
unifiée alliant dimensionnement solaire, gestion de projets et intelligence
artificielle. La méthode Kanban, pilotée sur Jira, assure un suivi visuel et
continu du flux de développement. Le chapitre suivant détaille la spécification
des besoins et l'architecture logicielle retenue.

\cleardoublepage
